\section{History of the T'au}

The xenos race known as the T'au, which inhabits a region of space near the Eastern Fringe of the galaxy, is young and expanding rapidly. Its weapons and technology are highly advanced. Although it is not yet two thousand years old, their nascent empire is spreading quickly through space and encountering the older races of the galaxy. This steady expansion has led to confrontations with the armies of the Imperium and those of other xenos races on many occasions, but strong in their advanced technology and driven by a common cause in which they place blind trust, the T'au extend their domain in the name of a Greater Good. This great idea sustains the T'au Empire: the individual must, in order to show a noble heart and to be just, set aside personal desires and work with others for the good of all. For the time being, this idea is respected by all. Compared with others, this empire is still limited and centred on a spherical cluster of stars, which allows the T'au to travel between them without resorting to Warp travel and the dangers that entails. Their domain includes several other xenos races that submitted of their own will or whose services were purchased through commercial agreements.

Unlike many xenos races encountered so far by Humanity, the T'au are not hostile, although they fiercely protect the territories they claim. The natural dynamism of the T'au drives them towards other occupied regions of the galaxy, which has inevitably created conflicts against humans as well as other peoples.

The T'au is humanoid in appearance, with two arms, two legs, and a single head, even though they have hooves in place of feet. Their skin is pale grey tending toward blue, although complexion and pigmentation can vary from one world to another; the inhabitants of Vior'la, for example, have much darker skin than natives of Bork'an. This tanned, dry epidermis is rough in texture, for it exudes practically no moisture. Smaller in stature than a human, the average T'au is slender, but their strength remains equivalent to that of an Imperial Guardsman, with a comparable tolerance for heat, cold, and pain. The face is flat at the level of the eyes, and it is thought that the spectrum of their visual field is slightly broader than that of a human in the infrared and ultraviolet. However, the absence of a pupil results in poor perception of distances and slower focusing. The T'au has no external olfactory organ; it is in fact located on the palate of the mouth, and at short range their sense of smell appears far more sensitive than that of a human. A T'au can therefore ``taste'' the air by guiding it with the tongue toward their sensory organs. Their life expectancy is generally shorter than that of a human, in particular those who have access to juvenats and other luxury medical treatments enjoyed by Imperial dignitaries, although a few rare individuals have found means of prolonging their lives well beyond the norm.

The appearance of a T'au reveals a great deal: caste, rank, home world, and status in society. The T'au commonly wears a single lock of hair adorned with ceremonial rings and trinkets that indicate rank and position, the most elaborate hair jewellery always denoting high responsibility. There are other, more specific ways of establishing a T'au's caste of origin: for example, Fire Caste warriors often bear deep scars, and it is not uncommon for the limbs of their leaders to be replaced by cybernetic implants. They most often come from the hot planets of their empire; their skin is usually darker than that of their kin and bears ritual tattoos and scarification. The artisans and labourers of the Earth Caste are stockier than other T'au, their hands more calloused, and they see life in a pragmatic way. This caste is less given to ornament, and few of these T'au indulge in any form of bodily decoration, which is considered frivolous and superfluous.

The Air Caste includes the pilots and messengers of the T'au, who spend most of their time in space. Their skin is generally paler and their bone structure lighter; some bones of their upper limbs are in fact hollow, which according to some of our scholars is a genetic remnant of a time when they were endowed with wings and could fly. When a member of the Air Caste is forced to remain on a planet's surface, they are extremely cautious in their movements, for long-term exposure to artificial gravity has reduced the resistance of the skeleton and made their bones fragile and brittle. The Water Caste is made up of merchants and diplomats, and it is generally the T'au of this caste who have travelled the most. Their features are softer and more expressive than those of the others. The respect these T'au have for other cultures makes them far more open to variety in dress and to the use of other dialects. It is thus common for a T'au merchant to prefer to negotiate while dressed in the same manner as their counterpart, and a diplomat often adopts the cultural manners of those with whom they deal when in their company. Of course, these practices always have an avowed purpose: the T'au way of life is never called into question and is held to be superior to all others.

The T'au of the last caste are the most mysterious. Their faces have the same typically flat features, but at the centre of their forehead stands a bony diamond. Its purpose, if it has one, remains an enigma, and despite several attempts the Inquisition has still been unable to recover the body of an Ethereal for thorough study. The gaps in knowledge of this more powerful and more important caste are a source of general consternation within the Ordo Xenos, and every scrap of information that can be gleaned is analysed with great attention.

Text from the Omnis Bibliotheca. More information at: \url{https://omnis-bibliotheca.com/index.php/Accueil}
